% BHU PYQ -- Transcribed & Corrected
% Paper: ESMJ21 / ESMN21: Mineral and Earth Material (Crystallography & Mineralogy)
% Exam: B.Sc. (Hons.) Semester II Examination, 2024-25 (NEP)
% Department: Department of Geology
% Paper Ref: CECONF/N/2024-25/26

\documentclass[11pt,a4paper]{article}
\usepackage[a4paper,top=2.5cm,bottom=2.5cm,left=2.5cm,right=2.5cm,headheight=14pt]{geometry}
\usepackage{amsmath,amssymb}
\usepackage[T1]{fontenc}
\usepackage[utf8]{inputenc}
\usepackage{lmodern,microtype,fancyhdr,enumitem,hyperref,lastpage,parskip}

\hypersetup{colorlinks=true,urlcolor=black,linkcolor=black,pdftitle={ESMJ21: Mineral & Earth Material -- BHU Sem II 2024-25 (NEP)}}
\pagestyle{fancy}\fancyhf{}
\fancyhead[L]{\small   \textit{Banaras Hindu University}}
\fancyhead[R]{\small   \textit{ESMJ21 / ESMN21: Mineral \& Earth Material}}
\fancyfoot[C]{\small Page  \thepage\ of \pageref{LastPage}}


\newlist{parts}{enumerate}{2}
\setlist[parts,1]{label=(\alph*),leftmargin=*,itemsep=4pt,topsep=2pt}
\setlist[parts,2]{label=(\roman*),leftmargin=*,itemsep=3pt,topsep=2pt}


\begin{document}


\begin{center}
  {\large\bfseries BANARAS HINDU UNIVERSITY}\\[2pt]
  {\normalsize B.Sc. (Hons.) Semester II Examination 2024-25 (NEP)}\\[6pt]
  \rule{0.8\linewidth}{0.5pt}\\[6pt]
  {\large\bfseries DEPARTMENT OF GEOLOGY}\\[4pt]
  {\large\bfseries Paper Code: ESMJ21 / ESMN21 --- Mineral and Earth Material}\\[4pt]
  {\normalsize Time: 2:30 Hours\hfill Maximum Marks: 70}\\[2pt]
  {\small Ref Code: CECONF/N/2024-25/26}
\end{center}

\rule{\linewidth}{0.4pt}\smallskip



\noindent   \textit{Instruction: Write your Roll No. at the top immediately on the receipt of this question paper.}\\[2pt]



\noindent   \textit{Note: Answer two questions each from Section--A and Section--B. Section--C is compulsory. Figures in the right-hand margin indicate marks.}
\bigskip

\section*{SECTION -- A}




\noindent   \textbf{Question 1} \hfill [14 Marks]\\
Provide a detailed explanation of the optical properties of anisotropic minerals observed under both Plane-Polarized Light (PPL) and Cross-Polarized Light (XPL), accompanied by appropriate diagrams.

\bigskip




\noindent   \textbf{Question 2} \hfill [14 Marks]\\
What are the diagnostic optical properties of plagioclase feldspar? Two orthoclase feldspar crystals (gemstone: 'Moonstone') were excavated from two pegmatite veins at different locations in an Indian craton. Do they possess the same interfacial angles? Justify your answer.

\bigskip




\noindent   \textbf{Question 3} \hfill [14 Marks]\\
``When an external force is applied to a mineral, the mineral responds.'' Describe this response in terms of any two physical properties of a mineral.

\bigskip




\noindent   \textbf{Question 4} \hfill [14 Marks]\\
Augite and Hornblende both possess prismatic crystal faces. Do they belong to the same crystallographic system? Differentiate between these two minerals based on their diagnostic physical and optical properties.

\bigskip
\section*{SECTION -- B}




\noindent   \textbf{Question 5} \hfill [14 Marks]\\
Explain / Answer in brief with suitable examples ($4  \times 3.5 = 14$ Marks):

\begin{parts}
  \item Why 5-fold, 7-fold, and 8-fold rotation axes are not possible in crystals?
  \item Explain the function of a reflecting goniometer and its uses.
  \item How would you differentiate between a closed form and an open form?
  \item Deduce the Miller indices from the following Weiss parameters:
  
\begin{romanlist}
    \item $1a, 2b, 3c$
    \item $2a, 4b, \infty c$
  \end{romanlist}
\end{parts}

\bigskip




\noindent   \textbf{Question 6} \hfill [14 Marks]\\
Answer the following with suitable diagrams and examples where necessary ($4  \times 3.5 = 14$ Marks):

\begin{parts}
  \item Differentiate between dispersion and diffusion of light.
  \item What is the function of the aperture diaphragm in a polarizing microscope?
  \item Why is Canada balsam commonly used as an adhesive in the preparation of thin sections in petrography?
  \item How is a Nicol prism constructed and what is its function in optical mineralogy?
\end{parts}

\bigskip




\noindent   \textbf{Question 7} \hfill [14 Marks]\\
Answer the following with suitable examples where necessary ($4  \times 3.5 = 14$ Marks):

\begin{parts}
  \item What is the basis for the classification of silicate minerals? Name the different classes of silicates.
  \item Why do diamond and graphite have the same chemical composition but different specific gravity?
  \item Why is the chemistry-based classification considered the most preferred classification scheme of minerals?
  \item Discuss the igneous and sedimentary processes of mineral formation.
\end{parts}

\bigskip
\section*{SECTION -- C (Compulsory)}




\noindent   \textbf{Question 8} \hfill [14 Marks]\\
Define / Answer the following in one word or sentence ($14  \times 1 = 14$ Marks):

\begin{parts}
  \item Minerals which are always black in crossed polars are called \rule{3cm}{0.15mm}.
  \item Which mineral shows polysynthetic twinning under the microscope?
  \item What is the role of the auxiliary condensing lens?
  \item How to differentiate between biotite and muscovite under the microscope?
  \item If you are given a task to distinguish a mineral and an amorphous substance, how will it be done?
  \item Name a non-ferromagnesian phyllosilicate mineral.
  \item What are the composition and hardness of a streak plate?
  \item What is the typical crystal form of garnet?
  \item If a mineral can scratch fluorite but not quartz, its hardness on the Mohs scale is most likely \rule{2cm}{0.15mm}.
  \item \rule{3cm}{0.15mm} is an imaginary three-dimensional framework representing positions of motifs in a crystal.
  \item A point where three or more crystal faces meet is called a \rule{3cm}{0.15mm}.
  \item Define pleochroic halo.
  \item Define birefringence in minerals.
  \item Give an example of an oxide mineral belonging to the spinel group.
\end{parts}

\end{document}
